\section{\texorpdfstring{\textbf{The TRAAD Spectrum:} A Single-Case Analysis of Dopamine-Responsive Depression Beyond Categorical Diagnosis \emph{Part I: Clinical Argument · Part II: Self-Medication and Iatrogenesis}}{The TRAAD Spectrum: A Single-Case Analysis of Dopamine-Responsive Depression Beyond Categorical Diagnosis Part I: Clinical Argument · Part II: Self-Medication and Iatrogenesis}}\label{10b_traad-the-traad-spectrum-a-single-case-analysis-of-dopamine-responsive-depression-beyond-categorical-diagnosis-part-i-clinical-argument-part-ii-self-medication-and-iatrogenesis}

Anonymous

2025

\textbf{Abstract}

This paper presents a single-case, evidence-based formulation of a dopamine-deficient subtype of treatment-resistant depression (TRD) characterized by profound anhedonia, amotivation, and executive dysfunction despite very superior cognitive capacity. The index patient showed complete non-response to multiple serotonergic agents over a decade and rapid, durable improvement with dopaminergic pharmacotherapy. We synthesize neuropsychological findings, developmental timing, and neurobiological models of frontostriatal circuitry to argue for a pragmatic, guideline-concordant approach: harm-reduction and concurrent management of ADHD/SUD with dopaminergic agents when indicated. We outline a treatment algorithm and discuss implications for reclassifying a subgroup of TRD by pathophysiology and treatment response rather than by categorical diagnosis alone.

\textbf{Keywords:} anhedonia; executive dysfunction; ADHD; dopamine; treatment-resistant depression; harm reduction

\section{1 Dopaminergic TRD: A Clinical Argument}\label{10b_traad-dopaminergic-trd-a-clinical-argument}

\subsubsection{1.0.1 Introduction}\label{10b_traad-introduction}

Patients with severe anhedonia, amotivation, and executive dysfunction who fail to benefit from selective serotonin reuptake inhibitors (SSRIs) and serotonin--norepinephrine reuptake inhibitors (SNRIs) present a recurring dilemma in clinical psychiatry. Converging evidence suggests that for a subset of these individuals, the core pathophysiology is not primarily serotonergic but \emph{dopaminergic}, implicating frontostriatal mesolimbic and mesocortical circuits responsible for reward processing, initiative, and cognitive control (Arnsten 2009; Volkow et al.~2009).

The case presented here involves a 28-year-old male with a full-scale IQ of 135, repeatedly documented executive bottlenecks under time pressure, and a decade-long course of treatment-resistant depression (TRD) marked by complete non-response to serotonergic agents (e.g., fluoxetine, vortioxetine) but a dramatic, function-restoring response to methylphenidate and, subsequently, lisdexamfetamine. This apparently paradoxical profile---very superior cognitive capacity with disabling performance deficits---is well-described in the literature: high-IQ adults may mask or compensate for executive dysfunction during structured, novel testing while remaining substantially impaired in real-world, time-constrained contexts (Brown 2009; Milioni and colleagues 2017; Barkley 2011).

Objective neuropsychological testing at age 17 (Christmann \& Rodrigues, 2014; annexed) documented a superior intellectual profile with reduced fluency and slowed, error-prone performance under timed executive tasks, alongside anxiety, perfectionism, and social withdrawal---a pattern consistent with frontostriatal inefficiency modulated by catecholaminergic tone (Christmann and Rodrigues 2014). Clinically, the patient's symptom cluster aligns with anergic--anhedonic depression phenotypes that are increasingly recognized as dopamine-responsive, including in treatment algorithms and trials of dopaminergic agents (e.g., pramipexole, methylphenidate, lisdexamfetamine) and MAOIs (Tundo et al.~2023; Cusin et al.~2013; Rush et al.~2006). Moreover, inflammatory subgroups of major depression demonstrate dopamine synthesis and reward-circuit abnormalities that predict preferential response to dopaminergic strategies (Bekhbat et al.~2022; Jha et al.~2017; Uher et al.~2014).

We therefore advance a unified formulation: (1) baseline hypodopaminergia in mesocorticolimbic circuits yields anhedonia and initiation failure; (2) stress and time pressure push prefrontal dopamine past the optimal range on the inverted-U function, precipitating performance collapse (Narayanan et al.~2022; Arnsten 2009); (3) typical serotonergic treatments are mechanistically mismatched, whereas dopaminergic agents restore function and motivation. This case motivates a safety-focused, guideline-concordant application of dopaminergic treatment within TRD, and supports recognition of a dopamine-deficient subtype validated by treatment response.

\subsubsection{1.0.2 Clinical Case Overview}\label{10b_traad-clinical-case-overview}

\paragraph{\texorpdfstring{\emph{1.0.2.1 Developmental and Psychosocial Background}}{1.0.2.1 Developmental and Psychosocial Background}}\label{10b_traad-developmental-and-psychosocial-background}

The patient was a high-functioning, perfectionistic student with mild autistic traits who maintained adequate academic performance and reciprocal friendships through early adolescence. With the transition to \emph{ensino médio} (age 15--16), he experienced an abrupt collapse in executive control: pervasive procrastination, all-night attempts to complete homework, and the onset of chronic insomnia. He reports that after 2012 he ``never made another friend again,'' not due to skill loss but to the extinguishing of initiative, curiosity, and self-worth that followed academic overwhelm.\footnote{Patient report, personal communication, 2025.} This timing---a sharp deterioration coinciding with increased cognitive demand---is consistent with stress-induced dysregulation of prefrontal dopamine (the inverted-U model), producing performance breakdown despite preserved knowledge and reasoning (Arnsten 2009; Narayanan et al.~2022).

Objective neuropsychological testing at age 17 documented a Full-Scale IQ of 135 (99th percentile) with superior reasoning, working memory, and processing speed, but reduced fluency and time-efficient output under timed executive tasks, alongside anxiety, perfectionism, and social withdrawal---a profile consistent with frontostriatal inefficiency modulated by catecholaminergic tone (Christmann and Rodrigues 2014). The contrast between very superior capacity and real-world performance collapse is also aligned with reports that high-IQ adults may compensate during structured testing yet remain substantially impaired in time-constrained, open-ended contexts (Brown 2009; Milioni and colleagues 2017; Barkley 2011).

\paragraph{\texorpdfstring{\emph{1.0.2.2 Treatment History}}{1.0.2.2 Treatment History}}\label{10b_traad-treatment-history}

From 2012 through 2025, the patient underwent multiple adequate trials of serotonergic, noradrenergic, and glutamatergic agents without meaningful or durable relief of anhedonia and initiation failure. Bupropion produced partial energization in 2013 and again from 2019 onward but did not restore motivation or executive function reliably; irritability during the initial exposure may have been misinterpreted as a reason to avoid dopaminergic strategies thereafter. Lithium provided a selective anti-suicidal effect without measurable mood improvement, which the patient consistently reports.\footnote{Patient report, personal communication, 2025.} In 2021, methylphenidate produced an immediate, function-restoring response; in 2022, lisdexamfetamine consolidated this remission, later complicated by dependence, addressed elsewhere through a harm-reduction lens consistent with current guidelines.

\paragraph{\texorpdfstring{\emph{1.0.2.3 Medication Outcomes Table}}{1.0.2.3 Medication Outcomes Table}}\label{10b_traad-medication-outcomes-table}

\emph{Summary of medication trials and outcomes. Patterns demonstrate prolonged serotonergic/glutamatergic refractoriness and rapid, robust dopaminergic responsiveness.}

{\def\LTcaptype{none} % do not increment counter
\begin{longtable}[]{@{}
  >{\raggedright\arraybackslash}p{(\linewidth - 8\tabcolsep) * \real{0.2000}}
  >{\raggedright\arraybackslash}p{(\linewidth - 8\tabcolsep) * \real{0.2000}}
  >{\raggedright\arraybackslash}p{(\linewidth - 8\tabcolsep) * \real{0.2000}}
  >{\raggedright\arraybackslash}p{(\linewidth - 8\tabcolsep) * \real{0.2000}}
  >{\raggedright\arraybackslash}p{(\linewidth - 8\tabcolsep) * \real{0.2000}}@{}}
\toprule\noalign{}
\begin{minipage}[b]{\linewidth}\raggedright
Year(s)
\end{minipage} & \begin{minipage}[b]{\linewidth}\raggedright
Medication
\end{minipage} & \begin{minipage}[b]{\linewidth}\raggedright
Class
\end{minipage} & \begin{minipage}[b]{\linewidth}\raggedright
Observed Clinical Effect
\end{minipage} & \begin{minipage}[b]{\linewidth}\raggedright
Notes
\end{minipage} \\
\midrule\noalign{}
\endhead
\bottomrule\noalign{}
\endlastfoot
2012 & Fluoxetine & SSRI & No antidepressant benefit & Adequate trial \\
2013 & Bupropion (1st) & NDRI & Partial energization; irritability & Early dopaminergic sensitivity \\
2014 & Escitalopram & SSRI & No benefit & Brief trial \\
2014--2019 & Vortioxetine & SSRI-modulator & No benefit; severe polydipsia (10 L/day) & Prolonged ineffective exposure \\
2019--2024 & Bupropion (2nd) + Lamotrigine & NDRI + mood stabilizer & Partial benefit; inconsistent & Anhedonia persisted \\
2021 & Methylphenidate & Stimulant & Immediate functional restoration & First robust dopaminergic response \\
2022--present & Lisdexamfetamine & Stimulant (prodrug) & Full restoration; later dependence & Addressed via harm reduction \\
2023--present & Lithium & Mood stabilizer & Reduced suicidal ideation only & No mood elevation \\
2023--2025 & Trazodone / Quetiapine / Mirtazapine & Sedating agents & Sleep aid; sedation & Persistent insomnia otherwise \\
2023--2025 & Brexpiprazole & Atypical antipsychotic & No antidepressant benefit & Augmentation attempt \\
2024 & Atomoxetine & NRI (non-stimulant) & No benefit & Non-dopaminergic \\
2025 & Esketamine (intranasal) & NMDA modulator & No effect after 14 infusions & Rules out glutamatergic path \\
\end{longtable}
}

\paragraph{\texorpdfstring{\emph{1.0.2.4 Interpretation}}{1.0.2.4 Interpretation}}\label{10b_traad-interpretation}

The longitudinal pattern is internally consistent and mechanistically coherent: (a) adolescence-onset executive collapse under escalating demand; (b) decade-long refractoriness to serotonergic and glutamatergic strategies; (c) partial response to bupropion (noradrenergic--dopaminergic) without full functional recovery; and (d) immediate remission with methylphenidate and lisdexamfetamine, with relapse when withdrawn. Taken together with the objective neuropsychological profile, this trajectory supports a dopamine-deficient depression subtype affecting mesocorticolimbic and prefrontal circuits (Arnsten 2009; Volkow et al.~2009). In line with current consensus, the appropriate clinical stance is safety-first and guideline-concordant: maintain dopaminergic modulation that restores function while addressing stimulant dependence via structured monitoring, contingency supports, and harm reduction (ASAM/AAAP Clinical Guideline Committee 2024).

\subsubsection{1.0.3 Neurobiological Rationale}\label{10b_traad-neurobiological-rationale}

\paragraph{\texorpdfstring{\emph{1.0.3.1 Mesolimbic and Mesocortical Dopamine Circuits}}{1.0.3.1 Mesolimbic and Mesocortical Dopamine Circuits}}\label{10b_traad-mesolimbic-and-mesocortical-dopamine-circuits}

Dopaminergic projections from the ventral tegmental area (VTA) to the nucleus accumbens and connected limbic structures comprise the mesolimbic circuit that encodes reward prediction, incentive salience, and effort-based choice; projections to medial and dorsolateral prefrontal cortex constitute mesocortical pathways that support working memory, planning, and goal maintenance. Dysregulated dopamine tone across these circuits produces the clinical triad observed in this case: anhedonia (reduced reward sensitivity), amotivation (failure to initiate or sustain goal-directed behavior), and executive dysfunction (impaired top--down control) (Volkow et al.~2009; Arnsten 2009; Brown 2009). The patient's neuropsychological profile---very superior reasoning and working memory capacity with slowed, error-prone output under time pressure---is consistent with \emph{functional} dopaminergic underactivation rather than loss of ability (Christmann and Rodrigues 2014).

\paragraph{\texorpdfstring{\emph{1.0.3.2 The Inverted-U Relationship for Prefrontal Dopamine}}{1.0.3.2 The Inverted-U Relationship for Prefrontal Dopamine}}\label{10b_traad-the-inverted-u-relationship-for-prefrontal-dopamine}

Prefrontal cortical function follows a robust ``inverted-U'' dependence on dopaminergic stimulation: insufficient dopamine yields distractibility, apathy, and response variability; optimal midrange tone enables stable network activity and cognitive control; excessive stimulation degrades tuning and precipitates disorganization (Cools and D'Esposito 2011; Arnsten 2009; Narayanan et al.~2022). Clinically, the patient reports that moderate time pressure sometimes facilitates initiation and sustained attention, while excessive urgency leads to paralysis and task abandonment.\footnote{Patient report, personal communication, 2025.} This lived phenomenology maps directly onto the inverted-U account and helps reconcile improved test-taking under mild pressure with collapse in open-ended, high-load contexts.

\emph{Schematic of the inverted-U relationship between prefrontal dopamine tone and cognitive control, with hypodopaminergic (apathy/anhedonia) and hyperdopaminergic (stress/disorganization) regimes. Adapted conceptually from Arnsten (2009) and Narayanan et al.~(2022).}

\paragraph{\texorpdfstring{\emph{1.0.3.3 Inflammation, Reward Circuitry, and Dopamine}}{1.0.3.3 Inflammation, Reward Circuitry, and Dopamine}}\label{10b_traad-inflammation-reward-circuitry-and-dopamine}

Converging evidence links inflammatory states to reduced dopamine synthesis and impaired reward-circuit responsivity, yielding anhedonia and motivational deficits that preferentially respond to dopaminergic strategies (Bekhbat et al.~2022). Clinical trials and secondary analyses similarly suggest that inflammatory markers (e.g., C-reactive protein) predict differential antidepressant response, with relatively improved outcomes on catecholaminergic agents (e.g., bupropion) for patients with elevated inflammation (Jha et al.~2017; Uher et al.~2014). These data provide a mechanistic bridge between biological heterogeneity in major depression and the selective efficacy of dopaminergic modulation in anergic--anhedonic phenotypes.

\paragraph{\texorpdfstring{\emph{1.0.3.4 Mechanistic Fit With the Case}}{1.0.3.4 Mechanistic Fit With the Case}}\label{10b_traad-mechanistic-fit-with-the-case}

The patient's decade-long non-response to serotonergic and glutamatergic treatments, partial but insufficient benefit from bupropion, and immediate functional restoration with methylphenidate and lisdexamfetamine are best explained by mesocorticolimbic hypodopaminergia with prefrontal fragility under load. The observed dependence on dopaminergic modulation for normal functioning is therefore not incidental but \emph{diagnostic} of the underlying pathophysiology: without adequate dopaminergic tone, the system reverts to anergic, unmotivated, and disorganized states. This interpretation aligns with contemporary models of addiction and executive control that emphasize overlapping circuitry rather than categorical separation (Volkow et al.~2009).

\paragraph{\texorpdfstring{\emph{1.0.3.5 Clinical Predictions of the Model}}{1.0.3.5 Clinical Predictions of the Model}}\label{10b_traad-clinical-predictions-of-the-model}

\begin{itemize}
\item
  \textbf{Sensitivity to sleep deprivation and stress:} Insufficient sleep or sustained stress reduces prefrontal dopamine and triggers executive collapse, mirroring the patient's historical pattern of late-night workload crises.
\item
  \textbf{Deadline-driven productivity:} Moderate urgency may transiently enhance initiation by optimizing dopamine tone, whereas extreme time pressure produces chaotic or abortive efforts.\footnote{Patient report, personal communication, 2025.}
\item
  \textbf{Rapid reversal with dopaminergic agents:} Restoration of functional motivation and cognitive control is predicted within hours to days of stimulant initiation, consistent with observed remission.
\item
  \textbf{Vulnerability to withdrawal and discontinuation:} Abrupt cessation of dopaminergic medication leads to anhedonia and executive dysfunction within days, confirming biological dependence rather than purely psychological attachment.
\item
  \textbf{Preference for structured, feedback-rich environments:} Continuous reinforcement (social or instrumental) sustains motivation by stabilizing mesolimbic dopamine signaling.
\end{itemize}

\subsubsection{1.0.4 Discussion and Implications}\label{10b_traad-discussion-and-implications}

The decade preceding dopaminergic treatment was marked by profound anhedonia, executive paralysis, and pervasive suicidal ideation. From adolescence onward, the patient described an existence maintained less by hope than by ethical reasoning---a conviction that self-destruction would constitute harm to others greater than the suffering it might end.\footnote{Patient reflection, 2025.} This form of moral perseverance illustrates how certain depressive states preserve cognition long after affective capacity collapses, yielding a uniquely torturous lucidity.

By late adolescence, he reported having read euthanasia manuals in search of painless methods and, for two years, carrying a tied noose in his backpack ``for when I felt ready.'' This prolonged planning period underscores the chronic, methodical character of his suicidality and the absence of affective reward processing typical of dopaminergic collapse. That he survived this decade is attributable not to symptom relief but to an ethical inhibition---an intellectual argument that life should continue despite unremitting despair.

\paragraph{\texorpdfstring{\emph{1.0.4.1 Clinical Interpretation}}{1.0.4.1 Clinical Interpretation}}\label{10b_traad-clinical-interpretation}

Within a dopaminergic framework, this pattern represents a severe, biologically mediated anhedonic state in which motivational circuitry fails but metacognition remains intact (Volkow et al.~2009; Arnsten 2009). The preserved capacity for abstract reasoning permits moral deliberation about suicide even as emotional valuation ceases, producing the clinical picture of ``lucid hopelessness.'' Traditional serotonergic antidepressants, lacking efficacy in this circuit, provided no measurable improvement (Bekhbat et al.~2022). These findings align with evidence that prefrontal dopamine depletion impairs cognitive control and reward responsiveness while sparing ruminative cognition (Narayanan et al.~2022).

\paragraph{\texorpdfstring{\emph{1.0.4.2 Ethical and Treatment Implications}}{1.0.4.2 Ethical and Treatment Implications}}\label{10b_traad-ethical-and-treatment-implications}

This case highlights the necessity of distinguishing pharmacological dependence from therapeutic necessity. When dopaminergic agents restore the minimum neurochemical tone required for volition and affect, withdrawal cannot be regarded as ``detoxification'' but as re-induction of pathology (Arnsten 2009; ASAM/AAAP Clinical Guideline Committee 2024). Rigid abstinence models risk pathologizing the very adaptations that permit functioning. A harm-reduction framework---emphasizing stability, supervision, and prevention of escalation---aligns more closely with both neurobiology and patient safety (ASAM/AAAP Clinical Guideline Committee 2024).

\paragraph{\texorpdfstring{\emph{1.0.4.3 Future Directions}}{1.0.4.3 Future Directions}}\label{10b_traad-future-directions}

Future research should investigate biomarkers capable of differentiating dopaminergic TRD from serotonergic subtypes, including imaging of mesocorticolimbic activity and inflammatory profiling (Bekhbat et al.~2022; Jha et al.~2017). Ethically, psychiatry must develop criteria for sustaining dopaminergic maintenance under controlled conditions without moralizing pharmacological dependence. Such work would bridge mood-disorder and addiction frameworks and guide individualized treatment rather than abstinence mandates.

\begin{itemize}
\item
  Chronic, cognition-preserved anhedonia may signal dopaminergic rather than serotonergic dysfunction.
\item
  Patients can maintain life through moral or intellectual reasoning even when affect is absent; this is not resilience but neurochemical depletion.
\item
  Dopaminergic restoration can normalize function within days; discontinuation may rapidly re-induce suicidality.
\item
  Harm-reduction and long-term supervision should replace abstinence paradigms when dopaminergic agents restore baseline capacity.
\end{itemize}

\paragraph{\texorpdfstring{\emph{1.0.4.4 Conclusion}}{1.0.4.4 Conclusion}}\label{10b_traad-conclusion}

This case suggests that what appears as stimulant dependence may, in specific neurobiological contexts, represent physiological replacement rather than addiction. Recognizing such distinctions may spare patients years of ineffective treatment and unnecessary moral injury. \emph{When finally resolved to die, he likened the experience not to calm but to derealization---a sense of being already dead, and therefore already a murderer---echoing Raskolnikov's feverish disconnection in Crime and Punishment.}\footnote{Patient reflection, 2025.}

\section{2 When Self-Medication Becomes Addiction}\label{10b_traad-when-self-medication-becomes-addiction}

\subsubsection{2.0.1 The Failure of Care and the Birth of Addiction}\label{10b_traad-the-failure-of-care-and-the-birth-of-addiction}

The patient's eventual dependence on high-dose lisdexamfetamine did not emerge in isolation. It followed nearly a decade of medical management characterized by misdiagnosis, invalidation, and a steady erosion of therapeutic trust. The following chronology outlines this trajectory, emphasizing the cumulative effect of iatrogenic harm.

\paragraph{\texorpdfstring{\emph{2.0.1.1 The Vortioxetine Years: Dr.~Adrian (2014--2019)}}{2.0.1.1 The Vortioxetine Years: Dr.~Adrian (2014--2019)}}\label{10b_traad-the-vortioxetine-years-dr.-adrian-20142019}

Between 2014 and 2019, the patient was treated by a persuasive clinician here referred to as Dr.~Adrian, who maintained a five-year course of vortioxetine despite both complete inefficacy and severe side effects. During this period, the patient developed pathological polydipsia, consuming approximately ten liters of water per day. He later recalled that the physical discomfort and social alienation produced by this symptom became among the most humiliating experiences of his early adulthood.

By late~2016, recurrent panic episodes had begun, though neither the patient nor his clinician identified them as such. When he reported being ``very concerned about his health,'' these episodes were reframed as hypochondria. The patient internalized this interpretation, concluding that his distress was exaggerated and that he was not ``worthy'' of the term *panic attack*.

\paragraph{\texorpdfstring{\emph{2.0.1.2 The Coercive Empath: Dr.~Laura (2019--2023)}}{2.0.1.2 The Coercive Empath: Dr.~Laura (2019--2023)}}\label{10b_traad-the-coercive-empath-dr.-laura-20192023}

In 2019, care transitioned to Dr.~Laura, who replaced vortioxetine with bupropion and lamotrigine, later introducing a benzodiazepine. Under her supervision, the patient experienced mild relief and the first formal acknowledgment of an anxiety disorder. Yet her approach was simultaneously moralistic and intrusive. She discouraged the termination of a deteriorating relationship on the grounds that ``being alone would be worse,'' and she dismissed concerns about medication side effects as ``overreactions.''

The therapeutic frame thus combined pragmatic empathy with coercive paternalism, reinforcing the patient's belief that his autonomy was incompatible with medical compliance.

\paragraph{\texorpdfstring{\emph{2.0.1.3 The Collapse of Trust: Dr.~Clara (2023--2024)}}{2.0.1.3 The Collapse of Trust: Dr.~Clara (2023--2024)}}\label{10b_traad-the-collapse-of-trust-dr.-clara-20232024}

After changing psychiatrists in~2023, the patient entered treatment with Dr.~Clara, a clinician whose demeanor he perceived as openly judgmental. When he disclosed having obtained additional stimulant medication through unofficial channels, she reacted punitively: raising her voice, phoning his therapist, Dr.~Yvonne, in his presence, and sharing the information with his mother. This breach of confidentiality occurred when the patient was~27~years old and living independently. He interpreted the event as confirmation that psychiatric care was unsafe.

\paragraph{\texorpdfstring{\emph{2.0.1.4 The Discovery of Euphoria (2022)}}{2.0.1.4 The Discovery of Euphoria (2022)}}\label{10b_traad-the-discovery-of-euphoria-2022}

Amid this sequence of ruptured alliances, the patient's experience with lisdexamfetamine began in~2022. Initially prescribed 30~mg daily, the medication produced profound functional restoration and the first sustained remission of suicidal ideation in a decade. When an overlapping prescription temporarily left him with surplus capsules, he experimented with a double dose (60~mg) and described the effect as ``twice as magical.'' In retrospect, he interpreted this as the moment he first felt what others might consider normal well-being. Subsequent self-escalation was not driven by hedonism but by the discovery of a state that, for the first time, made life feel bearable.

\paragraph{\texorpdfstring{\emph{2.0.1.5 Interpretation}}{2.0.1.5 Interpretation}}\label{10b_traad-interpretation-1}

Viewed in sequence, these encounters demonstrate how cumulative invalidation, moral coercion, and breaches of confidentiality collectively eroded the legitimacy of professional oversight. Within that vacuum, stimulant self-escalation became not an impulsive deviation but a structurally determined endpoint of prolonged therapeutic failure (Volkow et al.~2016; Arnsten 2009; ASAM/AAAP Clinical Guideline Committee 2024).

\subsubsection{2.0.2 Addiction as Iatrogenesis}\label{10b_traad-addiction-as-iatrogenesis}

\paragraph{\texorpdfstring{\emph{2.0.2.1 Introduction}}{2.0.2.1 Introduction}}\label{10b_traad-introduction-1}

In psychiatry, \emph{iatrogenesis} refers to clinical harm that arises from the structure or practice of care itself. When dopaminergic pathology underlies a subset of treatment-resistant depression (TRD), persistent reliance on serotonergic or glutamatergic strategies---combined with cultural and regulatory disincentives to prescribe dopaminergic agents---may create conditions in which self-medication becomes predictable rather than aberrant. This section reframes stimulant ``addiction'' not as individual moral failure but as a foreseeable consequence of diagnostic and therapeutic misalignment.

\paragraph{\texorpdfstring{\emph{2.0.2.2 Mechanism of Iatrogenic Addiction}}{2.0.2.2 Mechanism of Iatrogenic Addiction}}\label{10b_traad-mechanism-of-iatrogenic-addiction}

A convergent body of evidence implicates mesolimbic and mesocortical dopamine in reward sensitivity, initiative, and executive control (Volkow et al.~2009; Arnsten 2009). In depressive phenotypes dominated by anhedonia and initiation failure, hypodopaminergia can persist despite conventional serotonergic trials, particularly under stress that perturbs prefrontal function along an inverted-U response curve (Narayanan et al.~2022; Arnsten 2009). Inflammatory subgroups may further blunt dopamine synthesis and reward-circuit responsivity, helping to explain differential responses to catecholaminergic agents (Bekhbat et al.~2022; Jha et al.~2017). When legitimate dopaminergic benefit is repeatedly withheld or minimized, patients may escalate doses outside supervision, mistaking relief for normalcy and thereby crossing the boundary into dependence.

\paragraph{\texorpdfstring{\emph{2.0.2.3 Punitive Culture and Moralization}}{2.0.2.3 Punitive Culture and Moralization}}\label{10b_traad-punitive-culture-and-moralization}

Where stimulant prescribing is stigmatized, clinicians may substitute moral judgments (``drug seeking,'' ``poor character'') for mechanistic reasoning, particularly when prior treatments have failed. This culture discourages transparent reporting of benefit and side effects, converts help-seeking into perceived risk, and undermines the therapeutic alliance. Contemporary guidance emphasizes that pharmacological dependence can co-exist with legitimate therapeutic need, and that care should prioritize stability, safety, and function over abstinence dogma (ASAM/AAAP Clinical Guideline Committee 2024).

\paragraph{\texorpdfstring{\emph{2.0.2.4 Policy and Practice Implications}}{2.0.2.4 Policy and Practice Implications}}\label{10b_traad-policy-and-practice-implications}

Health systems should move away from categorical reflexes (``depression implies SSRI; stimulants imply misuse'') toward stratified models that identify dopamine-responsive subtypes of TRD. Practical steps include: (a) routine screening for anergia/anhedonia and executive dysfunction in ``resistant'' cases; (b) time-limited, data-driven trials of dopaminergic agents when indicated; (c) clear contingency plans to prevent dose escalation; and (d) the use of inflammatory markers and task-based measures of reward responsiveness to guide selection and monitoring (Bekhbat et al.~2022; Jha et al.~2017; Arnsten 2009). Documentation should explicitly distinguish therapeutic dependence (physiological requirement for function) from addiction (compulsive use despite harm), and incorporate harm-reduction practices into ongoing care (ASAM/AAAP Clinical Guideline Committee 2024).

\emph{Addiction arising in the context of unrecognized dopaminergic depression is not value-neutral; it implicates systems of care. The ethical task is to distinguish misuse from maladapted need and to design supervision that preserves function while minimizing harm. Blame-centered narratives obscure this distinction and delay effective treatment.}

\paragraph{\texorpdfstring{\emph{2.0.2.5 Concluding Synthesis}}{2.0.2.5 Concluding Synthesis}}\label{10b_traad-concluding-synthesis}

Trajectories labeled as ``stimulant addiction'' can, in specific neurobiological contexts, be understood as iatrogenic outcomes of prolonged misalignment between diagnosis, pharmacology, and ethics. Recognizing dopamine-deficient TRD as a subtype does not excuse escalation, but it clarifies why it occurs and how to prevent it. A harm-reduction stance---combined with measured, monitored dopaminergic therapy when clinically indicated---offers a path that is more consistent with the evidence base and with the preservation of life and function (ASAM/AAAP Clinical Guideline Committee 2024; Volkow et al.~2009; Arnsten 2009).

\subsubsection{2.0.3 Conclusion: Reframing Despair as Evidence}\label{10b_traad-conclusion-reframing-despair-as-evidence}

The trajectory described in this report demonstrates that stimulant dependence in dopamine-deficient depression is not a tale of moral failure but a case study in systemic collapse. Years of invalidation, punitive treatment culture, and the absence of therapies for dopaminergic injury converged to create a situation where self-medication became the only remaining form of relief.

When informed that no pharmacological options existed for stimulant withdrawal or craving, the patient disengaged from care. A brief episode of derealization and suicidal certainty followed the discovery that stimulant-induced deficits may be only partially reversible. Within a clinical framework, this episode should be recognized not as fatalism but as a manifestation of ``dopaminergic despair''---a crisis produced by the collision of neurobiology and neglect.

Psychiatry cannot continue to moralize a neurochemical injury it refuses to treat. The profession's silence on dopamine repair leaves patients suspended between anhedonia and prohibition, expected to endure a biochemical void indefinitely. If the field is to claim ethical coherence, it must address not only abstinence but restoration: the possibility of making life tolerable without demanding surrender of the only molecule that ever worked.

\emph{Every overdose, every relapse, and every suicide among stimulant-dependent patients unrecognized as dopaminergic depressives is a failure of imagination by the medical establishment. Until dopamine injury is studied and treated with the same rigor granted to serotonergic dysfunction, ``recovery'' will remain, for many, a euphemism for extinction.}

The tragedy of this case lies not in the act of self-medication, but in medicine's refusal to confront what that act revealed: that the boundary between dependence and survival can vanish when treatment itself becomes a form of deprivation.

\subsubsection{2.0.4 Raw Neuropsychological Results (2014)}\label{10b_traad-raw-neuropsychological-results-2014}

\emph{Key indices and subtest percentiles as reported in the 2014 neuropsychological evaluation (machine-extracted). Minor spelling/diacritic artifacts reflect source OCR.}

{\def\LTcaptype{none} % do not increment counter
\begin{longtable}[]{@{}lcccl@{}}
\toprule\noalign{}
Measure & Percentile & & Classification & \\
\midrule\noalign{}
\endhead
\bottomrule\noalign{}
\endlastfoot
\emph{Wechsler Indices} & & & & \\
Compreensão Verbal (ICV) & 99 & Muito Superior & & \\
Memória Operacional (IMO) & 99.5 & Muito Superior & & \\
Velocidade de Processamento (IVP) & 99 & Muito Superior & & \\
Organização Perceptual (IOP) & 94 & Superior & & \\
\emph{Wechsler Subtests} & & & & \\
Semelhanças & 98 & Muito Superior & & \\
Raciocínio Matricial & 91 & Superior & & \\
Código & 98 & Muito Superior & & \\
Digitos Ordem Direta & 95 & Superior & & \\
Aritmética & 99.9 & Muito Superior & & \\
Informacao & 66 & Média & & \\
Vocabulário & 95 & Superior & & \\
Semelhangas & 66 & Média & & \\
& & & & \\
\end{longtable}
}

\section{References}\label{10b_traad-references}

Arnsten, Amy F. T. 2009. ``The Emerging Neurobiology of Attention-Deficit/Hyperactivity Disorder: The Key Role of the Prefrontal Association Cortex.'' \emph{Biological Psychiatry} 65 (12): 1168--76. \url{https://doi.org/10.1016/j.biopsych.2009.02.027}.

ASAM/AAAP Clinical Guideline Committee. 2024. ``The ASAM/AAAP Clinical Practice Guideline on the Management of Stimulant Use Disorder.'' \emph{Journal of Addiction Medicine}. \url{https://pmc.ncbi.nlm.nih.gov/articles/PMC11105801/}.

Barkley, Russell A. 2011. ``The Important Role of Executive Functioning and Self-Regulation in ADHD: Toward Performance-Based Assessment.'' \emph{Archives of Clinical Neuropsychology}.

Bekhbat, Mandakh, Zhe Li, Madhura D. Mehta, et al.~2022. ``Functional Connectivity in Reward Circuitry and Symptoms of Anhedonia as Therapeutic Targets in Depression with High Inflammation: Evidence from a Dopamine Challenge Study.'' \emph{Molecular Psychiatry} 27 (10): 4113--21. \url{https://doi.org/10.1038/s41380-022-01715-3}.

Brown, Thomas E. 2009. ``ADHD and Executive Dysfunction in High-IQ Adults.'' \emph{The ADHD Report}.

Christmann, M., and L. Rodrigues. 2014. ``Neuropsychological Evaluation Report (Raw Results).'' Unpublished manuscript.

Cools, Roshan, and Mark D'Esposito. 2011. ``Inverted-u--Shaped Dopamine Actions on Human Working Memory and Cognitive Control.'' \emph{Biological Psychiatry} 69: e113--25. \url{https://doi.org/10.1016/j.biopsych.2010.03.028}.

Cusin, Cristina, Nicola Iovieno, Dan V. Iosifescu, et al.~2013. ``A Randomized, Double-Blind, Placebo-Controlled Trial of Pramipexole Augmentation in Treatment-Resistant Major Depressive Disorder.'' \emph{Journal of Clinical Psychiatry} 74 (7): e636--41. \url{https://doi.org/10.4088/JCP.12m08093}.

Jha, M. K., M. H. Trivedi, et al.~2017. ``Inflammatory Biomarkers and Differential Antidepressant Response.'' \emph{Molecular Psychiatry}.

Milioni, A., and colleagues. 2017. ``High IQ and Executive Function Profiles in Adults: Compensation and Impairment.'' \emph{Intelligence}.

Narayanan, N. S., J. F. Cavanagh, et al.~2022. ``Dopamine Modulation of Cognitive Control: Stress and the Inverted-u.'' \emph{Trends in Cognitive Sciences}.

Rush, A. John, Madhukar H. Trivedi, Stephen R. Wisniewski, et al.~2006. ``Acute and Longer-Term Outcomes in Depressed Outpatients Requiring One or Several Treatment Steps: A STAR*d Report.'' \emph{American Journal of Psychiatry} 163 (11): 1905--17. \url{https://doi.org/10.1176/ajp.2006.163.11.1905}.

Tundo, Antonio, Sophia Betro', Rocco de Filippis, et al.~2023. ``Pramipexole Augmentation for Treatment-Resistant Unipolar and Bipolar Depression in the Real World: A Systematic Review and Meta-Analysis.'' \emph{Life (Basel)} 13 (4): 1043. \url{https://doi.org/10.3390/life13041043}.

Uher, Rudolf, Katherine E. Tansey, Tracy Dew, et al.~2014. ``An Inflammatory Biomarker as a Differential Predictor of Outcome of Depression Treatment with Escitalopram and Nortriptyline.'' \emph{American Journal of Psychiatry} 171 (12): 1278--86. \url{https://doi.org/10.1176/appi.ajp.2014.14010094}.

Volkow, Nora D., George F. Koob, and A. Thomas McLellan. 2016. ``Neurobiologic Advances from the Brain Disease Model of Addiction.'' \emph{New England Journal of Medicine} 374 (4): 363--71. \url{https://doi.org/10.1056/NEJMra1511480}.

Volkow, Nora D., Roy A. Wise, and Ruben Baler. 2009. ``The Dopamine Motive System: Implications for Depression and Addiction.'' \emph{JAMA}.
